\section{First attempt: predicting price spikes (FuelGuard 1.0)}\label{app:v1}

Our first design asked machine learning to predict the \emph{price} directly. It is reported here because
its failure motivated the final product, and because the guidelines ask what the product cannot do.

\paragraph{Set-up.} Daily decisions on NYMEX diesel trading days, January 2008 -- October 2026 (4\,701
days), 21 raw series (diesel, Brent, WTI and gas futures for two maturities; diesel and Brent implied
volatility; copper, gold, 10-year Treasury futures; VIX and VIX3M; S\&P~500 and energy indices; US and euro
inflation expectations; ECB EUR/USD) and 57 engineered features. A decision at the close of $t$ is executed
at the settlement of $t+1$. The \emph{market cost} $M_t$ is the average diesel price in EUR/litre over
$t+2,\dots,t+22$; the \emph{lockable price} $P_t$ is a strip of first- and second-month futures settled at
$t+1$, each day priced with the contract that will be front month on that day (generic contracts roll on
the first trading day of the month; comparing different delivery months would create a spurious gain from
locking, a mistake our very first analysis made). Targets: the premium $R_t=M_t/P_t-1$ (regression) and a
\emph{spike} $\1\{R_t>5\%\}$ (classification, 19\% of days). Models: $L_2$-logistic and ridge, random forest,
gradient boosting, in the same purged expanding walk-forward (2012--2026). These $M_t$ and $P_t$ are reused
by the FuelGuard~2.0 backtest.

\paragraph{Results.} The premium was not predictable: the best out-of-sample $R^2$ against the futures price
was 0.5\%, no Diebold--Mariano test was significant, and cross-validation always chose the strongest ridge
penalty, i.e.\ a constant. Spikes were predictable only modestly (Table~\ref{tab:v1}): AUC 0.59--0.64 for the
57-feature models, positive Brier skill only for the random forest, while ranking days by diesel implied
volatility alone gave an AUC of 0.71. Skill collapsed in 2023 and 2026, when spikes started from calm
markets. In the hedging backtest a simple rule ``lock when implied volatility is above its median''
beat every classifier. The lesson: the predictable part of the problem is \emph{volatility}, not direction,
which led to FuelGuard~2.0.

\begin{table}[H]
\centering\small
\fittable{\tabdir/classification_oos.tex}
\caption{FuelGuard 1.0, spike classification out of sample 2012--2026 (3\,725 days). The one-feature
implied-volatility logistic was added after a first run. Last row: days ranked by implied volatility, no model.}
\label{tab:v1}
\end{table}
